\section{A Color-Coded R1CS-to-QAP Example}
\label{app:r1cs-qap-example}

This appendix gives a complete worked example of the transformation
\[
    \text{algebraic circuit}\longrightarrow \text{R1CS}\longrightarrow \text{QAP}.
\]
The example is intentionally color coded. The colors are not decoration: they are a reading device for tracking how one object is represented in the next layer. This is the arithmetization route underlying QAP-based systems such as GGPR, Pinocchio, and Groth16~\cite{GennaroGentryParnoRaykova2013,ParnoHowellGentryRaykova2013,Groth2016}. We stop at the QAP divisibility relation and do not construct a Groth16 proof.

\subsection{Color legend: wires, rows, and basis polynomials}

There are two independent color systems.

\begin{tcolorbox}
\textbf{How to read the colors.}
\begin{itemize}[leftmargin=2em]
    \item \textbf{Wire colors track columns.} The same wire has the same color in the algebraic circuit, in the assignment vector, in the R1CS matrix column, and in the QAP column polynomial.
    \item \textbf{Constraint colors track rows.} The same constraint color is used for the algebraic constraint $\mathcal C_i$, the R1CS row $i$, the constraint point $r_i$, and the Lagrange basis polynomial $L_i(X)$.
    \item In the QAP formula
    \[
        A_j(X)=\sum_{i=1}^{m} A_{i,j} L_i(X),
    \]
    the column index $j$ is the wire being encoded, while the basis polynomial $L_i(X)$ remembers which R1CS row supplied the coefficient.
\end{itemize}
\end{tcolorbox}

The wire-color legend is
\[
\begin{array}{c|ccccccc}
\text{wire} & \constwire{1} & \xone{x_1} & \xtwo{x_2} & \wone{w_1} & \vone{v_1} & \vtwo{v_2} & \ywire{y} \\
\hline
\text{role} & \text{constant} & \text{input} & \text{input} & \text{witness} & \text{intermediate} & \text{intermediate} & \text{output}
\end{array}
\]
The constraint-color legend is
\[
    \rowone{\mathcal C_1\leftrightarrow r_1=1\leftrightarrow L_1(X)},
    \qquad
    \rowtwo{\mathcal C_2\leftrightarrow r_2=2\leftrightarrow L_2(X)},
    \qquad
    \rowthree{\mathcal C_3\leftrightarrow r_3=3\leftrightarrow L_3(X)}.
\]
Thus a term such as $\rowtwo{L_2(X)}$ inside $A_{\xtwo{x_2}}(X)$ should be read as: ``the $\xtwo{x_2}$ column of the $A$ matrix has a nonzero coefficient in row $\rowtwo{\mathcal C_2}$.''

\subsection{The algebraic circuit}

Work over a field $\F$ of characteristic not equal to $2$ or $3$; for concrete numerical checks, one may read all values in $\F_{97}$. Consider the circuit
\[
    \ywire{y}=(\xone{x_1}+\xtwo{x_2})(\xtwo{x_2}+\wone{w_1}).
\]
Flatten the computation by introducing two intermediate wires:
\[
    \rowone{\mathcal C_1:}\qquad
    (\xone{x_1}+\xtwo{x_2})\cdot \constwire{1}=\vone{v_1},
\]
\[
    \rowtwo{\mathcal C_2:}\qquad
    (\xtwo{x_2}+\wone{w_1})\cdot \constwire{1}=\vtwo{v_2},
\]
\[
    \rowthree{\mathcal C_3:}\qquad
    \vone{v_1}\cdot\vtwo{v_2}=\ywire{y}.
\]
These three equations already have the R1CS shape
\[
    (\text{linear form})\cdot(\text{linear form})=(\text{linear form}).
\]
The concrete assignment used below is
\[
    \xone{x_1}=5,
    \qquad
    \xtwo{x_2}=6,
    \qquad
    \wone{w_1}=1,
\]
so that
\[
    \vone{v_1}=11,
    \qquad
    \vtwo{v_2}=7,
    \qquad
    \ywire{y}=77.
\]
The full assignment vector is
\[
    z=(\constwire{1},\xone{x_1},\xtwo{x_2},\wone{w_1},\vone{v_1},\vtwo{v_2},\ywire{y})
      =(1,5,6,1,11,7,77).
\]

\subsection{R1CS matrices with row and column labels}

Using the column order
\[
    (\constwire{1},\xone{x_1},\xtwo{x_2},\wone{w_1},\vone{v_1},\vtwo{v_2},\ywire{y}),
\]
the $A$ matrix encodes the left linear form in each constraint:
\[
\renewcommand{\arraystretch}{1.15}
\begin{array}{c|ccccccc}
A & \constwire{1} & \xone{x_1} & \xtwo{x_2} & \wone{w_1} & \vone{v_1} & \vtwo{v_2} & \ywire{y} \\
\hline
\rowone{\mathcal C_1} & 0 & \rowone{1} & \rowone{1} & 0 & 0 & 0 & 0 \\
\rowtwo{\mathcal C_2} & 0 & 0 & \rowtwo{1} & \rowtwo{1} & 0 & 0 & 0 \\
\rowthree{\mathcal C_3} & 0 & 0 & 0 & 0 & \rowthree{1} & 0 & 0
\end{array}
\]
Read the entry $A_{2,w_1}=\rowtwo{1}$ as follows: it is in the $\wone{w_1}$ column, so it is the coefficient of $\wone{w_1}$; it is in the $\rowtwo{\mathcal C_2}$ row, so it belongs to the second constraint.

The $B$ matrix encodes the right linear form in each constraint:
\[
\renewcommand{\arraystretch}{1.15}
\begin{array}{c|ccccccc}
B & \constwire{1} & \xone{x_1} & \xtwo{x_2} & \wone{w_1} & \vone{v_1} & \vtwo{v_2} & \ywire{y} \\
\hline
\rowone{\mathcal C_1} & \rowone{1} & 0 & 0 & 0 & 0 & 0 & 0 \\
\rowtwo{\mathcal C_2} & \rowtwo{1} & 0 & 0 & 0 & 0 & 0 & 0 \\
\rowthree{\mathcal C_3} & 0 & 0 & 0 & 0 & 0 & \rowthree{1} & 0
\end{array}
\]
The first two constraints multiply by the constant wire $\constwire{1}$; the third constraint multiplies by $\vtwo{v_2}$.

The $C$ matrix encodes the output linear form in each constraint:
\[
\renewcommand{\arraystretch}{1.15}
\begin{array}{c|ccccccc}
C & \constwire{1} & \xone{x_1} & \xtwo{x_2} & \wone{w_1} & \vone{v_1} & \vtwo{v_2} & \ywire{y} \\
\hline
\rowone{\mathcal C_1} & 0 & 0 & 0 & 0 & \rowone{1} & 0 & 0 \\
\rowtwo{\mathcal C_2} & 0 & 0 & 0 & 0 & 0 & \rowtwo{1} & 0 \\
\rowthree{\mathcal C_3} & 0 & 0 & 0 & 0 & 0 & 0 & \rowthree{1}
\end{array}
\]
The row colors show which equation a coefficient belongs to; the column colors show which wire the coefficient multiplies.

For the assignment $z=(1,5,6,1,11,7,77)$, the three rows check as follows:
\[
    \rowone{(A_1\cdot z)(B_1\cdot z)}=(5+6)\cdot1=11=\rowone{C_1\cdot z},
\]
\[
    \rowtwo{(A_2\cdot z)(B_2\cdot z)}=(6+1)\cdot1=7=\rowtwo{C_2\cdot z},
\]
\[
    \rowthree{(A_3\cdot z)(B_3\cdot z)}=11\cdot7=77=\rowthree{C_3\cdot z}.
\]
Thus this assignment satisfies the R1CS.

\subsection{From R1CS entries to QAP basis terms}

Choose the constraint domain
\[
    H=\{\rowone{r_1=1},\rowtwo{r_2=2},\rowthree{r_3=3}\}.
\]
The vanishing polynomial is
\[
    Z_H(X)=(X-1)(X-2)(X-3)=X^3-6X^2+11X-6.
\]
The Lagrange basis polynomials for $H$ are
\[
    \rowone{L_1(X)}=\rowone{\frac{(X-2)(X-3)}{(1-2)(1-3)}}=\rowone{\frac{(X-2)(X-3)}{2}},
\]
\[
    \rowtwo{L_2(X)}=\rowtwo{\frac{(X-1)(X-3)}{(2-1)(2-3)}}=\rowtwo{-(X-1)(X-3)},
\]
\[
    \rowthree{L_3(X)}=\rowthree{\frac{(X-1)(X-2)}{(3-1)(3-2)}}=\rowthree{\frac{(X-1)(X-2)}{2}}.
\]
They satisfy $L_i(r_j)=1$ if $i=j$ and $0$ otherwise.

The column-interpolation rule is
\[
\boxed{
    A_j(X)=\sum_{i=1}^{3} A_{i,j}\,L_i(X),
    \qquad
    B_j(X)=\sum_{i=1}^{3} B_{i,j}\,L_i(X),
    \qquad
    C_j(X)=\sum_{i=1}^{3} C_{i,j}\,L_i(X).
}
\]
Hence an R1CS number becomes a QAP coefficient: the entry $A_{i,j}$ becomes the coefficient of the row-colored basis polynomial $L_i(X)$ inside the wire-colored column polynomial $A_j(X)$.

For example, in the $A$ matrix the nonzero entries become:
\[
\begin{array}{c|c|c}
\text{R1CS entry} & \text{source} & \text{QAP contribution} \\
\hline
A_{1,x_1}=\rowone{1} & \xone{x_1}\text{ appears in }\rowone{\mathcal C_1} & \rowone{1\cdot L_1(X)} \\
A_{1,x_2}=\rowone{1} & \xtwo{x_2}\text{ appears in }\rowone{\mathcal C_1} & \rowone{1\cdot L_1(X)} \\
A_{2,x_2}=\rowtwo{1} & \xtwo{x_2}\text{ appears in }\rowtwo{\mathcal C_2} & \rowtwo{1\cdot L_2(X)} \\
A_{2,w_1}=\rowtwo{1} & \wone{w_1}\text{ appears in }\rowtwo{\mathcal C_2} & \rowtwo{1\cdot L_2(X)} \\
A_{3,v_1}=\rowthree{1} & \vone{v_1}\text{ appears in }\rowthree{\mathcal C_3} & \rowthree{1\cdot L_3(X)}
\end{array}
\]
This table is the key bridge. It shows exactly how a colored number in the R1CS matrix turns into a colored basis term in the QAP.

\subsection{Constructing the QAP column polynomials}

Now collect all contributions for each wire column.

For the $A$ matrix:
\[
\begin{array}{c|c}
\text{wire column} & A_j(X) \\
\hline
\xone{x_1} & \rowone{L_1(X)} \\
\xtwo{x_2} & \rowone{L_1(X)}+\rowtwo{L_2(X)} \\
\wone{w_1} & \rowtwo{L_2(X)} \\
\vone{v_1} & \rowthree{L_3(X)}
\end{array}
\]
All other $A_j(X)$ are zero. The term $\rowtwo{L_2(X)}$ inside $A_{\xtwo{x_2}}(X)$ is precisely the QAP image of the R1CS entry $A_{2,x_2}=1$.

For the $B$ matrix:
\[
\begin{array}{c|c}
\text{wire column} & B_j(X) \\
\hline
\constwire{1} & \rowone{L_1(X)}+\rowtwo{L_2(X)} \\
\vtwo{v_2} & \rowthree{L_3(X)}
\end{array}
\]
The constant wire appears in the first two $B$-linear forms, while $\vtwo{v_2}$ appears in the third.

For the $C$ matrix:
\[
\begin{array}{c|c}
\text{wire column} & C_j(X) \\
\hline
\vone{v_1} & \rowone{L_1(X)} \\
\vtwo{v_2} & \rowtwo{L_2(X)} \\
\ywire{y} & \rowthree{L_3(X)}
\end{array}
\]
For instance, $C_{\ywire{y}}(X)=\rowthree{L_3(X)}$ says that $\ywire{y}$ is the output wire of constraint $\rowthree{\mathcal C_3}$.

\subsection{Combining columns with the assignment}

The assignment values turn the column polynomials into three aggregate polynomials:
\[
    A_z(X)=\sum_j z_jA_j(X),
    \qquad
    B_z(X)=\sum_j z_jB_j(X),
    \qquad
    C_z(X)=\sum_j z_jC_j(X).
\]
Using the colored columns above,
\[
\begin{aligned}
A_z(X)
&=\xone{5}\,\rowone{L_1(X)}
 +\xtwo{6}\bigl(\rowone{L_1(X)}+\rowtwo{L_2(X)}\bigr)
 +\wone{1}\,\rowtwo{L_2(X)}
 +\vone{11}\,\rowthree{L_3(X)}\\
&=\rowone{11L_1(X)}+\rowtwo{7L_2(X)}+\rowthree{11L_3(X)}\\
&=4X^2-16X+23,
\end{aligned}
\]
\[
\begin{aligned}
B_z(X)
&=\constwire{1}\bigl(\rowone{L_1(X)}+\rowtwo{L_2(X)}\bigr)
 +\vtwo{7}\,\rowthree{L_3(X)}\\
&=\rowone{L_1(X)}+\rowtwo{L_2(X)}+\rowthree{7L_3(X)}\\
&=3X^2-9X+7,
\end{aligned}
\]
\[
\begin{aligned}
C_z(X)
&=\vone{11}\,\rowone{L_1(X)}+\vtwo{7}\,\rowtwo{L_2(X)}+\ywire{77}\,\rowthree{L_3(X)}\\
&=\rowone{11L_1(X)}+\rowtwo{7L_2(X)}+\rowthree{77L_3(X)}\\
&=37X^2-115X+89.
\end{aligned}
\]
The color of each $L_i$ term tells which R1CS row generated that contribution. Evaluating at the row-colored points reproduces the row values:
\[
\begin{array}{c|ccc}
X & A_z(X) & B_z(X) & C_z(X) \\
\hline
\rowone{1} & \rowone{11} & \rowone{1} & \rowone{11} \\
\rowtwo{2} & \rowtwo{7} & \rowtwo{1} & \rowtwo{7} \\
\rowthree{3} & \rowthree{11} & \rowthree{7} & \rowthree{77}
\end{array}
\]
Thus at each constraint point,
\[
    A_z(r_i)B_z(r_i)=C_z(r_i).
\]

\subsection{Checking the QAP divisibility condition}

Define the QAP numerator
\[
    P_z(X)=A_z(X)B_z(X)-C_z(X).
\]
Using the explicit polynomials above,
\[
\begin{aligned}
P_z(X)
&=(4X^2-16X+23)(3X^2-9X+7)-(37X^2-115X+89)\\
&=12X^4-84X^3+204X^2-204X+72.
\end{aligned}
\]
Direct evaluation gives
\[
    P_z(\rowone{1})=P_z(\rowtwo{2})=P_z(\rowthree{3})=0.
\]
Therefore $P_z$ is divisible by
\[
    Z_H(X)=X^3-6X^2+11X-6.
\]
Indeed,
\[
    P_z(X)=(12X-12)Z_H(X),
\]
so the quotient polynomial is
\[
    h_z(X)=12X-12.
\]
The final QAP relation for this example is
\[
\boxed{
    A_z(X)B_z(X)-C_z(X)=h_z(X)Z_H(X).
}
\]
The cross-representation trace is now visible:
\[
\begin{array}{c}
\text{colored algebraic constraint }\mathcal C_i\\
\Updownarrow\\
\text{colored R1CS row }i\\
\Updownarrow\\
\text{colored QAP basis polynomial }L_i(X).
\end{array}
\]
This is exactly the algebraic statement that the next Groth16 layer will commit to and check using pairings.

